\section{Introduction}\label{sec:introduction}

\subsection{What is Operations Research?}\label{subsec:introduction--what-is-operations-research}

Let's start with a definition of Operations Research (OR).

\highspace
\begin{definitionbox}[: Operations Research (OR)]
    \definition{Operations Research (OR)} is a branch of mathematics that uses \textbf{mathematical models} and \textbf{quantitative methods} to analyze complex decision-making problems and find optimal or near-optimal solutions.
\end{definitionbox}

\highspace
So \hl{Operations Research helps us make better decisions.} Suppose we have a real-world problem such as:
\begin{itemize}
    \item \emph{Which machine should execute each job?}
    \item \emph{How should a company allocate limited resources?}
    \item \emph{What is the shortest route from $A$ to $B$?}
    \item \emph{How many workers should be scheduled?}
    \item \emph{Where should a facility be built?}
\end{itemize}
In all these cases, there are many possible decisions. Some are infeasible, some are feasible but bad, and some are better than others. OR tries to make this decision process \textbf{systematic and quantitative}.

\highspace
In general, the idea is simple:
\begin{enumerate}
    \item Real-world problem.
    \item Formulate a mathematical representation.
    \item Analysis.
    \item Decision-making to find the best decision.
\end{enumerate}

\begin{examplebox}[: Manufacturing Decision Problem]
    Imagine a company deciding how much of three products to manufacture. The real-world description may contain things like: limited workforce, limited machine hours, profit per product, production times.

    \highspace
    \hl{OR translates those facts into mathematics.} We might introduce variables such as:
    \begin{equation*}
        x_1=\text{number of units of product 1},
    \end{equation*}
    \begin{equation*}
        x_2=\text{number of units of product 2},
    \end{equation*}
    And then construct an objective such as:
    \begin{equation*}
        \max \; 5x_1+8x_2
    \end{equation*}
    Subject to resource constraints.
\end{examplebox}

\highspace
From the above example, we can see that OR is not simply ``\emph{doing calculations}''. A large pat of OR is deciding \textbf{how to represent the decision problem mathematically in the first place}. To study decision problems, we have three major quantitative approaches:
\begin{itemize}
    \item \important{Optimization} (\emph{what should we do?}): we have decision variables, constraints, and an objective function, and we want the best feasible solution:
        \begin{equation*}
            \min \; \text{or} \; \max \; f(x) \quad \text{s.t. } x \in X
        \end{equation*}
        For example, minimize transportation cost while satisfying demand.

    \item \important{Game Theory} (\emph{what should we do when others are also making strategic decisions?}): we have \textbf{multiple decision-makers} whose choices interact. The outcome for one player depends on what the others do. For example, two competing companies choosing prices.

    \item \important{Simulation} (\emph{what will probably happen if we do this?}): we build a model of how a system behaves and then experiment with it. For example, simulate customers arriving at a call center to estimate waiting times under different staffing levels.
\end{itemize}
In these notes, we will focus on \textbf{optimization}.

\highspace
\begin{flushleft}
    \textcolor{Green3}{\faIcon{balance-scale} \textbf{Optimal vs Near-Optimal solutions}}
\end{flushleft}
In the previous definition, we mentioned that OR finds \textbf{optimal or near-optimal solutions}. So the natural question is: ``\emph{why not always optimal?}''. If we have an optimization problem:
\begin{equation*}
    \min_{x\in X} f(x)
\end{equation*}
An \textbf{optimal solution} $x^*$ satisfies:
\begin{equation*}
    f(x^*)\leq f(x) \qquad \forall x\in X
\end{equation*}
It is literally the best feasible solution. But some \hl{real optimization problems are computationally extremely difficult.} If the set of alternatives is enormous, finding the true optimum may \textbf{take too much time}. Then we may instead accept a solution that is:
\begin{itemize}
    \item Feasible;
    \item Computationally obtainable;
    \item Very close to the best known solution.
\end{itemize}
That is a \textbf{near-optimal solution}. Later in the notes this distinction will become much more important, especially when we discuss combinatorial and integer optimization.

\highspace
\begin{flushleft}
    \textcolor{Green3}{\faIcon{book} \textbf{OR as an interdisciplinary field}}
\end{flushleft}
Operations Research sits at the intersection of several areas: Applied Mathematics, Computer Science, Economics, and Industrial Engineering. Each contributes something different:
\begin{itemize}
    \item \textbf{Applied mathematics}. Mathematics gives us the language used to describe the problem. For example:
        \begin{equation*}
            \min f(x) \quad \text{s.t. } g_i(x)\leq 0
        \end{equation*}
        It also gives us the theoretical tools to determine whether solutions exist, whether they are optimal, and what properties they have.

    \item \textbf{Computer science}. A mathematical model is not enough. Suppose a problem has $2^{100}$ possible solutions. Writing the problem mathematically is easy; checking all solutions is impossible in practice. We therefore need algorithms to find good solutions efficiently.

        This is where computer science enters:
        \begin{equation*}
            \text{Model} \longrightarrow \text{Algorithm} \longrightarrow \text{Computer implementation}
        \end{equation*}
        These notes will repeatedly follow exactly this pattern. For example, later we will see shortest-path algorithms, minimum spanning-tree algorithms, maximum-flow algorithms, Simplex, and Branch-and-Bound.

    \item \textbf{Economics}. Many decision problems involve quantities such as: profit, cost, scarce resources, prices, trade-offs. So economic reasoning often determines what the objective function should represent. For instance:
        \begin{equation*}
            \max \text{ profit} \quad \text{or} \quad \min \text{ operating cost}
        \end{equation*}
        Later, LP duality will even give an explicit economic interpretation to certain mathematical quantities.

    \item \textbf{Industrial Engineering}. Industrial engineering provides many of the real problems that OR was designed to solve: scheduling, logistics, transportation, production planning, inventory, facility location, and supply chains.
\end{itemize}
So OR is not mathematics developed in isolation. It is \textbf{mathematics designed around real decision problems}.

\highspace
\begin{takeawaysbox}
    \begin{itemize}
        \item \textbf{Operations Research (OR)} uses mathematical models and quantitative methods to support complex decision-making and find \textbf{optimal or near-optimal solutions}.

        \item Its overall goal is to \textbf{help make better decisions} by translating a real-world problem into a mathematical model that can be analyzed and solved.

        \item The main quantitative approaches introduced in OR are:
            \begin{itemize}
                \item \textbf{Optimization}: find the best feasible decision;
                \item \textbf{Game Theory}: study interacting decisions of multiple decision-makers;
                \item \textbf{Simulation}: evaluate how a system behaves under alternative decisions.
            \end{itemize}

        \item A \textbf{near-optimal solution} may be accepted when finding the true optimum is computationally too expensive.

        \item OR is interdisciplinary, combining \textbf{applied mathematics, computer science, economics, and industrial engineering}.
    \end{itemize}
\end{takeawaysbox}